\scanpage{158}\mspage{32}
Ces propriétés s'explicitent ainsi :

\uline{M1} Propriétés de relèvement-prolongement
\[
C\mathrel{\unsure{\longleftarrow|\longleftrightarrow}}TF\quad\text{et}\quad TC\longleftrightarrow F
\]
\uline{M2} \unsure{Propriété|Propriétés} de factorisation de toute flèche $f$ de $M$ en $f=pi$, avec, au choix, $i\in C$, $p\in TF$ ou $i\in TC$, $p\in F$.

\[
M'\qquad TC\subset C\cap W,\quad TF\subset F\cap W
\]
\deleted{\uline{Ex} $W=TF\circ TC$ \quad(axiome d'exactitude)}

C'est donc une partie des axiomes d'une str. de Quillen faible \deleted{\uline{\unsure{close|stricte}}} pour $W$ ($\overline{\underline{\mathrm{XIII}}}$ p.27 déf. 1) -- Il manque les axiomes de stabilité M3,M4, et les axiomes de ``propriété'' (propriety) $M''$, savoir $F\subset\mathrm{Fib}_W$, $C\subset\mathrm{Cof}_W$.

\uline{Proposition 4.3.} Soient $M$ une catégorie, $(C,TF)$ et $(TC,F)$ deux \inserted{pré-}couples de Quillen dans $M$, satisfaisant
\[
TC\subset C,\quad TF\subset F
\tag*{(4.3.1.)}
\]
et soit $W$ une partie de $\mathrm{Fl}(M)$, dont on ne suppose pour l'instant que 4.6. b)\textsuperscript{*}, et
% slanted
\margin{* On utilise seulement que si $\unsure{X|x}\xrightarrow f Y\xrightarrow{\unsure{g|2}} Z$, et si $gf\in W$, alors $f\in W\Longleftrightarrow g\in W$. On n'a pas besoin que $W$ contienne les iso, ni qu'il soit stable par composition.}
\[
W\supset TC,TF,
\tag*{(4.3.2)}
\]
donc \deleted{$\unsure{F|T}$} on a
\[
TC\subset C\cap W,\quad TF\subset F\cap W.
\tag*{(4.3.3)}
\]
\deleted{(c'est à dire \unsure{de|les} propriétés de \unsure{saturation} pour $W$).}

Considérons les conditions de \deleted{factorisation} \inserted{prolongement\unsure{,|-} relèvement} suivantes a),a'), et les conditions de factorisation b),b'),\unsure{c)|b)},d) (en se rappelant que $\widetilde{TC}=F^*$, $\overline{TF}=C_*$, en vertu de prop. 4.2 Cor.1 appliqué aux \inserted{pré}couples de Quillen $(TC,F)$, $(C,TF)$) :
\[
\begin{aligned}
\text{a)}\quad&\deleted{\unsure{f|?}}\ C\cap W\subset\widetilde{TC}
&&\text{(donc }TC\subset C\cap W\subset\widetilde{TC}\text{)}\\
\text{a')}\quad&F\cap W\subset\overline{TF}
&&\text{(donc }TF\subset F\cap W\subset\overline{TF}\text{)}\\
\text{b)}\quad&W\subset\overset{\deleted{\unsure{\text{---}|\sim}}}{TF}\circ\widetilde{TC}\\
\text{b')}\quad&W\subset\overline{TF}\circ\overset{\deleted{\unsure{\sim|\text{---}}}}{TC}\\
\text{c)}\quad&W\subset\overline{TF}\circ\widetilde{TC}\\
\text{d)}\quad&W\subset TF\circ TC
\end{aligned}
\]
(donc $\overset{\inserted{\text{par }(4.3.2.)}}{W=TF\circ TC}$ si $W$ est stable par composition
% slanted
\margin{\uline{NB} \deleted{\unsure{comme|on}} a) signifie \unsure{juste} que $(C\cap W,F)$ est encore un précouple de Quillen, resp. \unsure{b)|a')} signifie que $(C,\unsure{F'|F}\cap W)$ en est un -- et alors cette paire de précouples de Quillen \unsure{satisfait} cette fois à d) i.e. $W\unsure{\subset|\in}TF'\circ TC'$.}
% archive sheet number: 157 (upper-left)
